\section{\ref{sec:code}장. 모스 부호}

이 장의 한계 추정은 정보 이론과 계산이다. 채널의 어디서 잡음이 생기는지, 뇌가 얼마나 빨리 작동하는지, 스파이크 하나가 얼마인지는 측정되었다. 피질이 실제로 어떤 부호를 쓰는지는 열린 문제로 남아 있다.

{\footnotesize
\setlength{\tabcolsep}{3pt}\renewcommand{\arraystretch}{1.15}
\begin{longtable}{>{\raggedright}p{0.5cm}>{\raggedright}p{4.0cm}>{\raggedright}p{5.6cm}>{\raggedright}p{2.3cm}>{\raggedright\arraybackslash}p{1.7cm}}
\toprule
No. & 명제 & 실험과 측정한 것 & 출처 & 상태 \\
\midrule\endfirsthead
\toprule
No. & 명제 & 실험과 측정한 것 & 출처 & 상태 \\
\midrule\endhead
1 & 피질 \nt{GLUT} 뉴런의 발화율은 초당 1개 미만에서 수십 개이며, 뉴런은 대부분의 시간을 아래쪽 경계 근처에서 작동한다 & 깨어 있는 쥐의 청각 피질에서 세포에 붙인 피펫으로 기록(뉴런 선택이 활동과 무관): 어느 순간이든 초당 $20$개가 넘는 발화율을 내는 뉴런은 $5\,\%$ 미만이다. 일부 뉴런의 배경 발화율은 초당 $0.01$보다 낮다. 발화율 분포는 로그 정규 분포이다 & \cite{Hromadka2008} & 측정됨 \\
2 & 스파이크 채널의 용량은 $R_{\max}\approx r\log_2\frac{e}{r\Delta t}$~\eqref{eq:spikeentropy}이고, 거의 전부가 스파이크 시각에 있다(명제~\ref{prop:capacity}) & 길이 $\Delta t$ 칸으로 된 이진 문자열의 엔트로피. 장의 수치: 초당 $r=10$, $\Delta t=1\ms$에서 스파이크당 $8$비트, $\Delta t=10\ms$에서 약 $5$비트, 스파이크 계수기는 초당 $2$비트 미만 & \cite{MacKay1952} & 이론 \\
3 & 감각 뉴런은 스파이크당 1비트에서 몇 비트를 전달하며, 가장 좋은 경우 한계의 절반까지 이른다 & 자극을 알고 있는 감각 뉴런의 스파이크로 자극을 복원; 책에 측정 요약이 있다 & \cite{Rieke1997} & 측정됨 \\
4 & 스파이크 발생기는 정확하다. 정확한 시각은 입력의 빠른 변화가 정한다 & 쥐 피질 절편의 뉴런: 빠르게 요동하는 반복 전류에는 스파이크가 $1\ms$보다 나은 정밀도로 반복되고, 일정한 전류에는 시각이 흩어진다 & \cite{Mainen1995} & 측정됨 \\
5 & 시냅스는 믿을 수 없다. 방출의 무작위성 때문에 스파이크의 절반 넘게가 수신자에 이르지 못한다 & 해마 피라미드 뉴런 입력의 최소 자극(절편): 스파이크의 절반 넘게가 응답을 일으키지 않는다. 자극 문턱의 요동, 축삭 가지점에서의 전도 실패, 낮은 실험 온도는 배제했다. 남는 것은 확률적 방출이다 & \cite{AllenStevens1994} & 측정됨 \\
6 & 살아 있는 피질의 스파이크 흐름은 무작위로 보인다. 간격은 푸아송 흐름처럼 흩어지고, 스파이크 수의 분산은 평균에 가깝다. “잡음”의 일부는 동물의 움직임에 관한 메시지이다 & 깨어 있는 원숭이의 시각 영역 V1과 MT 기록: 간격의 변동 계수는 약 $1$이고, 스파이크 수의 분산 대 평균 비는 무작위 과정과 같다. 독립적인 작은 입력을 합하는 모델은 훨씬 작은 산포를 준다. 생쥐 시각 피질에서는 변동의 상당 부분이 동물 자신의 움직임을 찍은 영상으로 예측된다 & \cite{Softky1993,Stringer2019} & 측정됨 \\
7 & 발화율 부호는 느리다. 오차는 $1/\sqrt{rT}$~\eqref{eq:rateerror}인데, 뇌는 $\sim150\ms$ 만에 그림을 알아본다 & 공식은 푸아송 통계로, 장에서 유도했다. 실험: 사람이 $20\ms$ 동안 보여 준 낯선 사진에 동물이 있는지 판단했다. “예”와 “아니오”에 대한 뇌 전위는 약 $150\ms$ 뒤에 갈라진다. 이 시간을 $10$--$20\ms$씩 열 개쯤의 단계로 나눈 것은 장의 추정이지 측정이 아니다 & \cite{Thorpe1996} & 측정됨 \\
8 & 무리에 대한 평균은 상관으로 제한된다. 오차는 기껏해야 $1/\sqrt c$배 줄어든다~\eqref{eq:correlated}(명제~\ref{prop:pooling}) & 원숭이 MT 영역에서 동시에 기록한 뉴런 쌍: 스파이크 수의 상관 계수는 평균 $0.12$이며, 이론적 분석은 이 정도 상관만으로도 무리에서 얻는 이득이 제한됨을 보여 준다. 이론: 한계를 정하는 것은 상관 가운데 신호를 닮은 부분뿐이다. 반대 입장의 모델: 뉴런 $50$--$100$개 무리의 발화율은 스파이크 간격 하나, $10$--$50\ms$ 안에 읽히며, 피질에서 개별 스파이크의 시각은 쓰이지 않는다 & \cite{Zohary1994,MorenoBote2014,ShadlenNewsome1998} & 측정됨 \\
9 & 수는 첫 스파이크 시각~\eqref{eq:latency}, 무리의 스파이크 순서, 국소 리듬의 위상에 적을 수 있다 & 도롱뇽 망막: 짧게 보여 준 그림의 공간 구조가 출력 뉴런 첫 스파이크의 상대 지연에 적히며, 대비와 거의 무관하다. 쥐 해마의 장소 세포: “자기” 장소를 지날 때 스파이크가 세타 리듬($7$--$12$\,Hz) 주기 안에서 점점 일찍 옮겨 가며, 이동 폭은 $100^\circ$에서 $355^\circ$이다. 피질이 스파이크 시각을 얼마나 쓰는지는 열린 문제이다(8행) & \cite{Gollisch2008,OKeefe1993} & 측정됨 \\
10 & 어떤 부호가 읽힐지는 수신자의 시간 상수가 정한다~\eqref{eq:reader}(명제~\ref{prop:reader}). 활성 피질의 뉴런은 동시성 검출기에 더 가깝다 & 식~\eqref{eq:reader}는 장에서 유도했다. in vivo 기록 총설: 활성 피질에서 막 컨덕턴스는 휴지 상태보다 몇 배 크고, 시간 상수는 그만큼 짧다. 피질이 바로 이런 식으로 부호를 전환한다는 것은 직접 보이지 않았다 & \cite{Destexhe2003} & 가설 \\
11 & 고갈 시냅스는 발화율의 변화, 본질적으로 $\Delta r/r$를 전달한다~\eqref{eq:depression}(그림~\ref{fig:depression}) & 피질 피라미드 뉴런의 쌍 기록: 고갈 속도는 방출 확률이 정한다. 고갈이 느리면 응답이 발화율을 반영하고, 빠르면 일치를 반영한다. 쥐 피질에서 측정한 고갈에 기초한 모델: 빠른 입력과 느린 입력에서 같은 상대 발화율 변화가 같은 응답을 준다. 즉 시냅스는 $\Delta r/r$를 전달한다. 수치 $1.7\to2.2$, $6.7$, $90\ms$는 장의 계산이다 & \cite{Tsodyks1997,Abbott1997depression} & 측정됨 \\
12 & 버스트는 “대시”이다. 버스트가 사라질 확률 $(1-p)^k$~\eqref{eq:burst}는 작고, 촉진이 이를 더 줄인다 & 총설: 믿을 수 없는 시냅스에서 단일 스파이크는 자주 사라지지만, 버스트는 촉진 덕분에 믿을 만하게 전달된다. 버스트는 시냅스의 장기 변화를 일으킨다. 뇌가 버스트를 별도의 기호로 읽는다는 것은 가설이다 & \cite{Lisman1997,AllenStevens1994} & 가설 \\
13 & 빠른 \nt{GABA} 뉴런은 리듬과 “구두점”을 정한다. 또 다른 부류는 억제하는 뉴런을 억제하여 관문을 연다(명제~\ref{prop:punctuation}) & 피질 \nt{GABA} 뉴런 부류들의 성질에 관한 총설. 광유전학: 빠른 \nt{GABA} 뉴런을 리듬 있게 흥분시키면 감마 리듬이 생기고, 자극에 대한 응답은 주기의 위상에 의존한다. 생쥐 시각 피질에서 PV 뉴런은 서로를, SST 뉴런은 나머지 모든 부류를, VIP 뉴런은 주로 SST를 억제한다. 깨어 있는 생쥐에서 VIP 뉴런을 흥분시키면 \nt{GLUT} 뉴런의 억제가 풀린다. VIP 뉴런은 강화 신호로 켜진다. 일반 규칙으로서의 역할 분담은 가설이다 & \cite{Tremblay2016,Cardin2009,Pfeffer2013,Pi2013} & 가설 \\
14 & 에너지는 평균 발화율을 초당 스파이크 1개 정도로 제한한다~\eqref{eq:energy} & 해부학적, 생리학적 자료로 본 설치류 회색질의 에너지 예산: 스파이크와 시냅스 전류가 신호 전달 에너지의 $47\,\%$와 $34\,\%$이다. 동시에 활성일 수 있는 뉴런은 $\sim15\,\%$를 넘지 못한다. 사람 피질에서는 동시에 강하게 활성일 수 있는 뉴런이 $\sim1\,\%$를 넘지 못한다. 스파이크당 ATP $\sim10^9$개와 신호 전달 전력 $\sim10$\,W는 이 계산에 기초한 장의 추정이다 & \cite{Attwell2001,Lennie2003} & 이론 \\
15 & 부호는 희소하며 줄당 비트 수가 최대가 되도록 맞추어져 있다. 피질은 정밀도를 에너지와 맞바꾼다(명제~\ref{prop:economy}) & 절약 부호 가설. 근거: 먹이를 제한한 생쥐의 시각 피질에서 AMPA 수용체 컨덕턴스와 시냅스의 ATP 소비가 줄고($-29\,\%$), 스파이크 발화율은 유지되며, 방위 조율은 $32\,\%$ 넓어진다 & \cite{Barlow1961,Padamsey2022} & 가설 \\
\bottomrule
\end{longtable}}
